% concatenated from kFU-group-GOR.tex
\documentclass[preprint,11pt]{elsarticle}


\usepackage{amsfonts, amsmath, amscd}
\usepackage[psamsfonts]{amssymb}

\usepackage{amssymb}

\usepackage{pb-diagram}

\usepackage[all,cmtip]{xy}

%% The amsthm package provides extended theorem environments
\usepackage{amsthm}

%\usepackage[inline]{showlabels}


\usepackage[usenames]{color}


\headheight=0in
\headsep = 0.51in
\topmargin=0in
\textheight=8.950in
\textwidth=6.5in
\oddsidemargin=-0.19in
\evensidemargin=-0.19in
%\textwidth=7.18in
%\oddsidemargin=-0.19in
%\evensidemargin=-0.19in
\parindent=0.2in


%\journal{Journal of Algebra}


\newtheorem{theorem}{Theorem}[section]
\newtheorem{lemma}[theorem]{Lemma}
\newtheorem{corollary}[theorem]{Corollary}
\newtheorem{question}[theorem]{Question}
%\newtheorem{remark}[theorem]{Remark}
\newtheorem{proposition}[theorem]{Proposition}
%\newtheorem{definition}[theorem]{Definition}
%\newtheorem{example}[theorem]{Example}
\newtheorem{fact}[theorem]{Fact}
\newtheorem{claim}[theorem]{Claim}
\newtheorem{notation}[theorem]{Notation}
\newtheorem{problem}[theorem]{Problem}
\newtheorem{conjecture}[theorem]{Conjecture}

\theoremstyle{definition}
\newtheorem{example}[theorem]{Example}
\newtheorem{definition}[theorem]{Definition}
\theoremstyle{remark}
\newtheorem{remark}[theorem]{Remark}




\numberwithin{equation}{section}
\numberwithin{theorem}{section}
%\numberwithin{example}{section}


%%%%%%%%%%%%%%%%%%%%%%%%%%%%%%%%%%%%%%%%%%%%%%%
%%%%%%%%%%%%%%%%%%%%%%%%%%%%%%%%%%%%%%%%%%%%%%%
%%%%%%%%%%%%%%%%%%%%%%%%%%%%%%%%%%%%%%%%%%%%%%%
%%%%%%%%%%%%%%%%%%%%%%%%%%%%%%%%%%%%%%%%%%%%%%%
%%%%%%%%%%%%%%%%%%%%%%%%%%%%%%%%%%%%%%%%%%%%%%%
%%%%%%%%%%%%%%%%%%%%%%%%%%%%%%%%%%%%%%%%%%%%%%%
%%%%%%%%%%%%%%%%%%%%%%%%%%%%%%%%%%%%%%%%%%%%%%%
%%%%%%%%%%%%%%%%%%%%%%%%%%%%%%%%%%%%%%%%%%%%%%%

\newcommand{\e}{\varepsilon}
\newcommand{\w}{\omega}
\newcommand{\onetwo}{\circ\bullet\!}
%\newcommand{\onetwo}{\mbox{\rm\textonehalf}\!}

%%%%%%%%%%%%%%%%%%%%%%%%%%%%
\newcommand{\NN}{\mathbb{N}}
\newcommand{\ZZ}{\mathbb{Z}}
\newcommand{\RR}{\mathbb{R}}
\newcommand{\VV}{\mathbb{V}}
\newcommand{\PP}{\mathbb{P}}
\newcommand{\IR}{\mathbb{R}}
\newcommand{\IN}{\mathbb N}
\newcommand{\II}{\mathbb{I}}
\newcommand{\Ss}{\mathbb{S}}
\newcommand{\ff}{\mathbb{F}}
\newcommand{\Q}{\mathbb{Q}}
\newcommand{\TT}{\mathbb{T}}
\newcommand{\ID}{\mathbb{D}}
\newcommand{\IF}{\mathbb{F}}
\newcommand{\IC}{\mathbb{C}}


%%%%%%%%%%%%%%%%%%%%%%%%%%%%
\newcommand{\CP}{\mathbf{cp}}
\newcommand{\dd}{(\mathbf{D})}
\newcommand{\MM}{\mathbf{M}}
\newcommand{\yyy}{\mathbf{y}}
\newcommand{\zzz}{\mathbf{z}}
\newcommand{\aaa}{\mathbf{a}}
\newcommand{\uuu}{\mathbf{u}}
\newcommand{\ball}{\mathbf{B}}
\newcommand{\sphere}{\mathbf{S}}
\newcommand{\eee}{\mathbf{e}}



%%%%%%%%%%%%%%%%%%%%%%%%%%%%
\newcommand{\SSS}{\mathfrak{S}}
\newcommand{\GG}{\mathfrak{G}}
\newcommand{\Ff}{\mathfrak{F}}
\newcommand{\nn}{\mathfrak{n}}
\newcommand{\Pp}{\mathfrak{P}}
\newcommand{\LLL}{\mathfrak{L}}
\newcommand{\RRR}{\mathfrak{R}}
\newcommand{\Tt}{\mathfrak{T}}
\newcommand{\pp}{\mathfrak{p}}
\newcommand{\con}{\mathfrak{c}}


%%%%%%%%%%%%%%%%%%%%%%%%%%%%
\newcommand{\TTT}{\mathcal{T}}
\newcommand{\FF}{\mathcal{F}}
\newcommand{\F}{\mathcal{F}}
\newcommand{\V}{\mathcal{V}}
%\newcommand{\U}{\mathcal{U}}
\newcommand{\W}{\mathcal{W}}
\newcommand{\BB}{\mathcal{B}}
\newcommand{\DD}{\mathcal{D}}
\newcommand{\UU}{\mathcal{U}}
\newcommand{\WW}{\mathcal W}
\newcommand{\EE}{{\mathcal E}}
\newcommand{\SC}{\mathcal{S}}
\newcommand{\KK}{\mathcal{K}}
\newcommand{\Nn}{\mathcal{N}}
\newcommand{\AAA}{\mathcal{A}}
\newcommand{\I}{\mathcal{I}}
\newcommand{\LL}{\mathcal{L}}
\newcommand{\M}{\mathcal{M}}
\newcommand{\A}{\mathcal{A}}
%\newcommand{\C}{\mathcal{C}}
\newcommand{\Gg}{\mathcal{G}}
\newcommand{\PPP}{\mathcal{P}}
\newcommand{\Se}{\mathcal{S}}
\newcommand{\Ope}{\mathcal{Op}}

%%%%%%%%%%%%%%%%%%%%%%%%%%%%%%%%%%%%
\newcommand{\supp}{\mathrm{supp}}
\newcommand{\conv}{\mathrm{conv}}
\newcommand{\acx}{\mathrm{acx}}
\newcommand{\const}{\mathrm{const}}
\newcommand{\pr}{\mathrm{pr}}
\newcommand{\cl}{\mathrm{cl}}
%\newcommand{\Lin}{\mathrm{Lin}}
\newcommand{\Ra}{\Rightarrow}
\newcommand{\La}{\Leftarrow}
\newcommand{\LRa}{\Leftrightarrow}
\newcommand{\cacx}{\overline{\mathrm{acx}}}
\newcommand{\Hom}{\mathrm{Hom}}
\newcommand{\rng}{\mathrm{rng}}
\newcommand{\id}{\mathrm{id}}


%%%%%%%%%%%%%%%%%%%%%%%%%%%%%%%%%%%
\newcommand{\Fi}{\mathsf{Fi}}
\newcommand{\Nuc}{\mathsf{Nuc}}
\newcommand{\OP}{\mathsf{Op}}
\newcommand{\Ll}{\mathsf{L}}
\newcommand{\HS}{\mathsf{HS}}
\newcommand{\Co}{\mathsf{Com}}
\newcommand{\co}{\mathsf{co}}
\newcommand{\Fin}{\mathsf{Fin}}
\newcommand{\Bo}{\mathsf{Bo}}
\newcommand{\FB}{\mathsf{FB}}
\newcommand{\Int}{\mathsf{Int}}
\newcommand{\diam}{\mathsf{diam}}
\newcommand{\dom}{\mathsf{dom}}
\newcommand{\Mac}{\mathsf{Mac}}
\newcommand{\sgn}{\mathsf{sign}}
\newcommand{\qc}{\mathsf{qc}}
\newcommand{\Id}{\mathsf{id}}
\newcommand{\Lin}{\mathsf{Lin}}
\newcommand{\Sch}{\mathsf{Sch}}
\newcommand{\LQC}{\mathsf{LQC}}
\newcommand{\MAPA}{\mathsf{MAPA}}
\newcommand{\Nor}{\mathsf{Nor}}
\newcommand{\Ban}{\mathsf{Ban}}
\newcommand{\ind}{\mathsf{ind}}
\newcommand{\proj}{\mathsf{proj}}
\newcommand{\IM}{\mathsf{Im}}
\newcommand{\spn}{\mathsf{span}}
%\newcommand{\csp}{\overline{\mathrm{span}}}
\newcommand{\cspn}{\overline{\mathsf{span}}}

\newcommand{\LCS}{\mathsf{LCS}}
\newcommand{\SL}{\mathsf{S}}
\newcommand{\SLCS}{\mathsf{SL}}

%%%%%%%%%%%%%%%%%%%%%%%%%%%%%%%%%%%%
\newcommand{\CC}{C_k}



\newcommand{\mus}{\pmb{\mu}_{\mathbf{s}}}
\newcommand{\nus}{\pmb{\nu}_{\mathbf{s}}}
\newcommand{\VVs}{\VV_{\mathbf{s}}}
\newcommand{\VVS}{\VV(\mathbf{s})}
\newcommand{\cantor}{2^{\omega}}

\newcommand{\SI}{\underrightarrow{\mbox{ s-$\mathsf{ind}$}}_n\,}
\newcommand{\SM}{{\setminus}}


%%%%%%%%%%%%%%%%%%%%%%%%%%%%%%%%%%%%

\newcommand{\weakstar}{\ensuremath{weak^*}}
%\newcommand{\Int}{\operatorname{int}}
%\newcommand{\diam}{\operatorname{diam}}
%\newcommand{\eps}{\varepsilon}
%\renewcommand{\phi}{\varphi}
\newcommand{\la}{\langle}
\newcommand{\ra}{\rangle}
\newcommand{\Op}{\mathcal{O}\/}
\newcommand{\uhr}{\upharpoonright}
\newcommand{\tcr}{\textcolor{red}}


%\newcommand{\diam}{\mathrm{diam}}

\newcommand{\GEN}{\mathsf{G_{E}}}
\newcommand{\SGEN}{\mathsf{G^s_{E}}}
\newcommand{\GNE}{\mathsf{G_{N}}}
\newcommand{\BM}{\mathsf{BM}}
\newcommand{\Sol}{\mathsf{Sol}}
\newcommand{\Gsa}{\widehat{G}_{\sigma^\ast}}




%%%%%%%%%%%%%%%%%% Notations
%%%%%%%%%%%%%%%%%%
%%%%%%%%%%%%%%%%%%
%%%%%%%%%%%%%%%%%%
%%%%%%%%%%%%%%%%%%

\newcommand{\ZN}{\ZZ^{(\NN)}}


%%%%%%%%%%%%%%%%%%
%%%%%%%%%%%%%%%%%%
%%%%%%%%%%%%%%%%%%
%%%%%%%%%%%%%%%%%%


\def\om{\omega}
\def\al{\alpha}
\def\cl#1{\overline{#1}}
\def\la{\lambda}
\def\es{\varnothing}
\def\Int{\operatorname{Int}}
\def\de{\delta}
\def\ep{\varepsilon}

\newcommand\restrB[2]{\ensuremath{\left.#1\right|_{#2}}}
\def\restr#1#2{\restrB{#1}{#2}}

%%%%%%%%%%%%%%%%%%
%%%%%%%%%%%%%%%%%%
%%%%%%%%%%%%%%%%%%
%%%%%%%%%%%%%%%%%%
%%%%%%%%%%%%%%%%%%

%%%%%%%%%%%%%%%%%%
%%%%%%%%%%%%%%%%%%
%%%%%%%%%%%%%%%%%%
%%%%%%%%%%%%%%%%%%
%%%%%%%%%%%%%%%%%%

\begin{document}

\begin{frontmatter}


\title{On $\kappa$-Fr\'{e}chet--Urysohn topological groups}


\author{Saak Gabriyelyan}%\tnotetext[label1]{The first named author was partially supported by Israel Science Foundation grant 1/12.}
\ead{saak@math.bgu.ac.il}
\address{Department of Mathematics, Ben-Gurion University of the Negev, Beer-Sheva, Israel}

\author{Alexander V. Osipov}%\tnotetext[label1]{The first named author was partially supported by Israel Science Foundation grant 1/12.}
\ead{OAB@list.ru}
\address{Krasovskii Institute of Mathematics and Mechanics, Ural Federal University, Yekaterinburg, Russia}


\author{Evgenii Reznichenko}%\tnotetext[label1]{The second named author was partially supported by Israel Science Foundation grant 1/12.}
\ead{erezn@inbox.ru}
\address{Department of Mathematics, Lomonosov Mosow State University, Moscow, Russia}


\begin{abstract}
We characterize $\kappa$-Fr\'{e}chet--Urysohn topological groups. Using this characterization we show that: (1) a hemicompact topological group is $\kappa$-Fr\'{e}chet--Urysohn iff it is locally compact, and (2) if $F$ is a closed metrizable subspace of a topological vector space (tvs) $E$ such that the quotient $E/F$ is a $\kappa$-Fr\'{e}chet--Urysohn space, then also $E$ is a $\kappa$-Fr\'{e}chet--Urysohn space. Consequently, the product of a $\kappa$-Fr\'{e}chet--Urysohn tvs and a metrizable tvs is a $\kappa$-Fr\'{e}chet--Urysohn space. Under Martin's Axiom, we construct a countable Boolean $\kappa$-Fr\'{e}chet--Urysohn group which is not a $k_\IR$-space.
\end{abstract}

\begin{keyword}
topological group \sep topological vector space \sep $\kappa$-Fr\'{e}chet--Urysohn  \sep $k_\IR$-space


\MSC[2020] 22A05 \sep 54A05  \sep   54C35 \sep 54D99

\end{keyword}

\end{frontmatter}



%%%%%%%%%%%%%%%%%%
%%%%%%%%%%%%%%%%%%
%%%%%%%%%%%%%%%%%%
%%%%%%%%%%%%%%%%%%
%%%%%%%%%%%%%%%%%%

\section{Introduction}

%%%%%%%%%%%%%%%%%%
%%%%%%%%%%%%%%%%%%
%%%%%%%%%%%%%%%%%%
%%%%%%%%%%%%%%%%%%
%%%%%%%%%%%%%%%%%%

The class of $\kappa$-Fr\'{e}chet--Urysohn spaces was introduced by Arhangel'skii. A Tychonoff space $X$ is called {\em $\kappa$-Fr\'{e}chet--Urysohn} if, for any open set $U\subseteq X$ and each point $x\in \overline{U}$, there exists a sequence $\{x_n\}_{n\in \w}\subseteq U$ that converges to~$x$. It should be mentioned that without using the term $\kappa$-Fr\'{e}chet--Urysohn, Mr\'{o}wka proved in \cite{Mrowka} that any product of first countable spaces is  $\kappa$-Fr\'{e}chet--Urysohn. $\kappa$-Fr\'{e}chet--Urysohn spaces are thoroughly studied by Liu and Ludwig \cite{LiL}, Sakai \cite{Sak2,Sakai}, Gabriyelyan \cite{Gabr-B1,Gabr-seq-Ascoli}, Gabriyelyan and Reznichenko \cite{GR-compact}, and Gabriyelyan, Osipov and Reznichenko \cite{GOR-Baire}.

In this article we continue the study of topological groups (in particular, topological vector spaces) which are $\kappa$-Fr\'{e}chet--Urysohn spaces. In Theorem \ref{t:kFU-group-charac} we characterize $\kappa$-Fr\'{e}chet--Urysohn topological groups. Using this characterization we show in Proposition \ref{p:kFU-hemicompact} that a hemicompact topological group $G$ is $\kappa$-Fr\'{e}chet--Urysohn if, and only if, $G$ is Baire if, and only if, $G$ is locally compact. It is known that the product of a Fr\'{e}chet--Urysohn space and a non-discrete metrizable space can be not $\kappa$-Fr\'{e}chet--Urysohn, see \cite{GR-compact}. However, in the realm of topological vector spaces the situation changes. Indeed, we show in Corollary \ref{c:prod-kFU-metr-tvs} that if $E$ is a $\kappa$-Fr\'{e}chet--Urysohn topological vector space and $F$ is a metrizable topological vector space, then the product $E\times F$ is a $\kappa$-Fr\'{e}chet--Urysohn space. Under Martin's Axiom, we construct a countable Boolean $\kappa$-Fr\'{e}chet--Urysohn group which is not a $k_\IR$-space, see Example \ref{exa:kFU-count-non-kR}.

%%%%%%%%%%%%%%%%%%
%%%%%%%%%%%%%%%%%%
%%%%%%%%%%%%%%%%%%
%%%%%%%%%%%%%%%%%%
%%%%%%%%%%%%%%%%%%

\section{Main results} \label{sec:kFU-charac}

%%%%%%%%%%%%%%%%%%
%%%%%%%%%%%%%%%%%%
%%%%%%%%%%%%%%%%%%
%%%%%%%%%%%%%%%%%%
%%%%%%%%%%%%%%%%%%

%All topological spaces in the article are assumed to be Tychonoff.

We start with the following characterization of topological groups which are $\kappa$-Fr\'{e}chet--Urysohn spaces, we shall call such groups by {\em  $\kappa$-Fr\'{e}chet--Urysohn groups}. The identity of a group $G$ is denoted by $e$.

Following  \cite{Osipov-23}, a family $\AAA=\{A_n\}_{n\in\w}$ of subsets of a space $X$ {\em weakly converges} to a point $x$ if for every neighborhood $W$ of $x$, there are a sequence $\{a_n\}_{n\in\w}$ and $m\in\w$ such that $a_n\in A_n$ $(n\in\w)$ and $a_n\in W$ for every $n>m$.

%Below we show that $\kappa$-Fr\'{e}chet--Urysohn groups can be characterized by a property similar to the property $(\alpha_4)$ of Arhangel'skii.
%The next theorem generalizes a characterization of  $\kappa$-Fr\'{e}chet--Urysohn  topological vector spaces obtained in \cite{GOR-Baire}.

\begin{theorem} \label{t:kFU-group-charac}
For a topological group $G$, the following assertions are equivalent:
\begin{enumerate}
\item[{\rm(i)}] for every sequence $\{\mathcal{V}_n: n\in \w\}$ of countable families $\mathcal{V}_n=\{V^i_n: i\in \w\}$ of open subsets of $G$ such that $\mathcal{V}_n$ weakly converges to $e$ for every $n\in\w$, there are strictly increasing sequences  $(i_k)$ and $(n_k)$ in $\w$ and a sequence $\{x_k\}_{k\in\w}$ such that $x_k\to e$ and $x_k\in V^{i_k}_{n_k}$ for every $k\in\w$;
\item[{\rm(ii)}] for every sequence $\{U_n\}_{n\in\w}$ of open subsets of $G$ such that $e\in \overline{U_n}$ for every $n\in\w$, there are a strictly increasing sequence $(n_k)\subseteq \w$ and a sequence $\{x_k\}_{k\in\w}$ in $G$ such that $x_k\in U_{n_k}$ $(k\in\w)$ and $x_k\to e$;
\item[{\rm(iii)}] $G$ is a $\kappa$-Fr\'{e}chet--Urysohn space.
\end{enumerate}
\end{theorem}

\begin{proof}
(i)$\Ra$(ii) For every $i\in \w$, we set $V^i_n:=U_n$. Now (ii) immediately follows from (i).
\smallskip

(ii)$\Ra$(iii) Let $U$ be an open subset of $G$ such that $e\in \overline{U}\SM U$. For every $n\in\w$, set $U_n:=U$. Then (iii) immediately follows from (ii).
\smallskip

(iii)$\Ra$(i) Let $\{\mathcal{V}_n: n\in \w\}$  be a family satisfying (i). Choose an arbitrary one-to-one sequence $\{p_n\}_{n\in\w}$ in $G\SM \{e\}$  converging to $e$ (such a sequence exists by the $\kappa$-Fr\'{e}chet--Urysohness applying to the open set $G\SM\{e\}$). For every $n\in\w$, set $p_n\cdot\mathcal{V}_n:=\{p_n V^i_n: i\in \w\}$. Choose a sequence $\{W_n: n\in \w\}$ of open neighborhoods of $e$ such that
\begin{equation} \label{equ:kFU-group-charac-1}
%W_{n+1}+W_{n+1}\subseteq W_n, \;\; (p_n+W_n) \cap (p_{n+1}+W_{n+1})=\emptyset \;\; \mbox{ and }\;\;
e\not\in \overline{p_n W_n}
\end{equation}
for every $n\in \w$. For each $0\leq n\leq i$, set
\begin{equation} \label{equ:kFU-group-charac-2}
\mathcal{O}^i_n=p_nV^i_n\cap p_nW_n.
\end{equation}
We claim that $e\in \overline{\bigcup \{\mathcal{O}^i_n: 0\leq n\leq i\}}$.  Indeed, let $U$ be a neighborhood of $e$. Choose a neighborhood $U_1$ of $e$ such that $U_1\cdot U_1\subseteq U$. Since $p_n\to e$, there is $m\in\w$ such that $p_n\in U_1$ for each $n\geq m$.  As $\mathcal{V}_m$ weakly converges to $e$, there are $j\geq m $ and a point $s_m^j\in V_m^j$ such that $s_m^j\in U_1\cap W_m$. Then $p_m\cdot s_m^j\in U_1\cdot U_1\subseteq U$ and $p_m\cdot s_m^j\in \mathcal{O}^j_m$.
\smallskip

Since $G$ is $\kappa$-Fr\'{e}chet--Urysohn, there is a sequence $(z_k)\subseteq \bigcup \{\mathcal{O}^i_n: 0\leq n\leq i\}$ converging to $e$. For every $k\in\w$, choose indices  $i_k$ and $n_k$ such that $z_k\in \mathcal{O}^{i_k}_{n_k}$. The sequence $(n_k)$ cannot be bounded because of (\ref{equ:kFU-group-charac-1}) and the inclusion $\mathcal{O}^{i_k}_{n_k} \subseteq p_{n_k}W_{n_k}$. Taking into account that $j_k\geq n_k$ it follows that also $(j_k)$ is unbounded. Passing to subsequences if needed, we can assume that $(j_k)$ and $(n_k)$ are strictly increasing.  For every $k\in\w$, set $x_k:= p_{n_k}^{-1}\cdot z_k $. Then, by (\ref{equ:kFU-group-charac-2}), $x_k\in V^{i_k}_{n_k}$ and, clearly, $x_k\to e$. Thus (i) is satisfied.
\end{proof}

%By Example 14.3 of \cite{kak}, every Fr\'{e}chet--Urysohn hemicompact abelian topological group is a locally compact Polish space. The natural counterpart of this result for $\kappa$-Fr\'{e}chet--Urysohness is the following assertion.

As an application of Theorem \ref{t:kFU-group-charac} we obtain a structure of hemicompact groups with the $\kappa$-Fr\'{e}chet--Urysohn property.

\begin{proposition} \label{p:kFU-hemicompact}
For a hemicompact topological group  $G$ the following assertions are equivalent:
\begin{enumerate}
\item[{\rm(i)}] $G$ is $\kappa$-Fr\'{e}chet--Urysohn;
\item[{\rm(ii)}]  $G$ is locally compact;
\item[{\rm(iii)}] $G$ is Baire.
\end{enumerate}
\end{proposition}

\begin{proof}
(i)$\Ra$(ii) Let $G$ be $\kappa$-Fr\'{e}chet--Urysohn, and suppose for a contradiction that $G$ is not locally compact. Take an increasing sequence $\{K_n\}_{n\in\w}$ of compact subsets of $G$ such that any compact subset of $G$ is contained in some $K_n$. We assume that $e\in K_0$. By our supposition all $K_n$ are not neighborhoods of $e$. For every $n\in\w$, set $U_n:=G\SM K_n$. Then $U_n$ is an open subset of $G$ such that $e\in \overline{U_n}\SM U_n$. Since  $G$ is $\kappa$-Fr\'{e}chet--Urysohn, (ii) of Theorem \ref{t:kFU-group-charac} implies that there are a strictly increasing sequence $(n_k)\subseteq \w$ and a sequence $\{x_k\}_{k\in\w}$ in $G$ such that $x_k\in U_{n_k}$ $(k\in\w)$ and $x_k\to e$. Since $K:=\{x_k\}_{k\in\w}\cup\{e\}$ is a compact subset of $G$, there is $m\in \w$ such that $K\subseteq K_m$. Take $k\in\w$ such that $n_k>m$. Then $x_k\in K_m\cap  U_{n_k}=K_m\cap (G\SM  K_{n_k})=\emptyset$. This is a contradiction.

(ii)$\Ra$(i) If $G$ is locally compact, then $G$ is $\kappa$-Fr\'{e}chet--Urysohn by Corollary 6.6 %\ref{c:lcg-kFU}
of \cite{GR-compact}.

(ii)$\Ra$(iii) is trivial.

(iii)$\Ra$(ii) If $G$ is Baire, then (see (i)), there is $n\in\w$ such that the compact set $K_n$ contains an open subset. Thus $G$ is locally compact.
\end{proof}
The condition of being $\kappa$-Fr\'{e}chet--Urysohn in Proposition \ref{p:kFU-hemicompact} cannot be replaced by the property of being a sequential space. Indeed, it is well-known that the free abelian group $A([0,\w])$ over the convergent sequence $[0,\w]$ is a sequential hemicompact group which is not locally compact.

Example 14.3 of \cite{kak} states that an abelian hemicompact topological group $G$ is Fr\'{e}chet--Urysohn if, and only if, it is a locally compact Polish group. Recall that a Tychonoff space $X$ is {\em Fr\'{e}chet--Urysohn} if, for any $A\subseteq X$ and each $x\in \overline{A}$, there exists a sequence $\{a_n: n\in \w\}\subseteq A$ that converges to $x$. Below we show that the condition of being abelian can be omitted.


\begin{corollary} \label{c:FU-hemicompact}
Let $G$ be a hemicompact topological group. Then $G$ is Fr\'{e}chet--Urysohn if, and only if, it is a locally compact Polish group.
\end{corollary}

\begin{proof}
Assume that $G$ is Fr\'{e}chet--Urysohn. Then, by Proposition \ref{p:kFU-hemicompact}, $G$ is locally compact. The group $G$ being Fr\'{e}chet--Urysohn has countable tightness. Recall that, by \cite{Ismail}, the character of any locally compact group coincides with its tightness. Therefore $G$ has countable character, and hence, by the Kakutani metrization theorem, $G$ is metrizable. As $G$ is hemicompact and metrizable, it is separable. Thus $G$ is  a locally compact Polish group. The converse assertion is clear.
\end{proof}

Proposition 14.4 of \cite{kak} states that if $F$ is a closed metrizable subspace of a tvs $E$ such that $E/F$ is Fr\'{e}chet--Urysohn, then the space $E$ is a Fr\'{e}chet--Urysohn space, too. Below we show that an analogous result holds also for $\kappa$-Fr\'{e}chet--Urysohn spaces.

\begin{proposition} \label{p:kFU-metriz-quot}
If $F$ is a closed metrizable subspace of a topological vector space $E$ such that the quotient $E/F$ is a $\kappa$-Fr\'{e}chet--Urysohn space, then  $E$ is also a $\kappa$-Fr\'{e}chet--Urysohn space.
\end{proposition}

\begin{proof}
Since $F$ is a metrizable subspace of $E$, there exists a decreasing sequence $\{V_n\}_{n\in\w}$ of balanced neighborhoods of zero in $E$ such that $V_{n+1}+V_{n+1}\subseteq V_n$ for all $n\in \w$ and the sequence $\{V_n\cap F\}_{n\in\w}$ is a basis of neighborhoods of zero in $F$. To prove that $E$ is
$\kappa$-Fr\'{e}chet--Urysohn, it is sufficient to show that if $U\subseteq E$ is an open set and $0\in \overline{U}\SM U$, then there exists a sequence in $U$ converging to $0$.

Let $p: E\rightarrow E/F$ be the quotient map. %If $0\in \overline{U\cap F}$, then the metrizability of $F$ immediately implies the existence of a desired sequence in $U$ converging to zero. If $0\not\in \overline{U\cap F}$, then we can replace $U$ by $U\SM \overline{U\cap F}$ and assume that $U\cap F=\emptyset$.
Note that $0\in \overline{V_n\cap U}$  for all $n\in \omega$. Hence $p(0)\in \overline{p(V_n\cap U)}$. Since $V_n\cap U$ is open and $p$ is an open map, the set $p(V_n\cap U)$ is open in $E/F$. Hence, by Theorem \ref{t:kFU-group-charac}(ii),  there are a sequence $(n_k)\subseteq \w$ and  a sequence $\{x_{k}\}_{k\in\w}\subseteq E$ such that
\[
p(x_{k})\in p(V_{n_k}\cap U) \; \mbox{ and } \; p(x_{k})\rightarrow p(0) \mbox{ as } k\to \infty.
\]
In particular, for every $k\in\w$, there are $z_k\in V_{n_k}\cap U$ and  $f_k\in F$ such that $x_k=z_k+f_k$. Replacing $x_k$ by $z_k$ if needed, we can assume that $x_k\in V_{n_k}\cap U$ for every $k\in\w$.


Fix a balanced neighborhood $W$ of zero in $E$. Choose $n_0\in \w$ such that $(V_{n_0}+V_{n_0})\cap F\subseteq W$. Taking into account that $p(W\cap V_{n_0})$ is an open neighborhood of zero in $E/F$, $W\cap V_{n_0}$ is balanced and $p(x_{k})\to 0$, we obtain that there is an $m>n_0$ such that
\[
\big[x_k+(W\cap V_{n_0})\big]\cap F\neq \emptyset \quad \mbox{ for all $k>m$.}
\]
Therefore, for every $k>m$, there exist $y_k\in W\cap V_{n_0}$ and  $u_k\in F$, such that $x_k+y_k=u_k$. Since $x_k\in V_{n_k}\subseteq V_k\subseteq V_{n_0}$ and $y_k\in V_{n_0}$, we deduce that $u_k=x_k+y_k\in (V_{n_0}+V_{n_0})\cap F\subseteq W$. Therefore $x_k=u_k-y_k\in W+W$. Since $W$ was arbitrary, this implies that $\{x_k\}_{k\in \w}$ is a sequence in $U$ that converges to $0$ in $E$. Thus $E$ is a $\kappa$-Fr\'{e}chet--Urysohn space.
\end{proof}

Michael showed (see \cite[Corollary~14.3]{kak}) that the product of a Fr\'{e}chet--Urysohn tvs and a metrizable tvs is Fr\'{e}chet--Urysohn. Proposition \ref{p:kFU-metriz-quot} shows that an analogous result holds true also for $\kappa$-Fr\'{e}chet--Urysohn tvs.

\begin{corollary} \label{c:prod-kFU-metr-tvs}
If $E$ is a $\kappa$-Fr\'{e}chet--Urysohn topological vector space and $F$ is a metrizable topological vector space, then the product $E\times F$ is a $\kappa$-Fr\'{e}chet--Urysohn space.
\end{corollary}
Note that in Corollary \ref{c:prod-kFU-metr-tvs}, the condition that $E$ and $F$ are topological vector spaces is essential. Indeed, in Theorem 3.7 of \cite{GR-compact} it is shown that the product of the Fr\'{e}chet--Urysohn fun $V(\w)$ with an arbitrary non-discrete space $X$ is not $\kappa$-Fr\'{e}chet--Urysohn.

Being motivated by the famous Malykhin problem (which asks whether there exists a countable non-metrizable Fr\'{e}chet--Urysohn group), the following problem arises naturally (note that this problem complements Problem 2.1 of \cite{GR-compact}).
\begin{problem} \label{prob:kFU-not-kR}
Does there exist a countable  $\kappa$-Fr\'{e}chet--Urysohn group which is not metrizable?
\end{problem}
Under Martin's Axiom, in Example \ref{exa:kFU-count-non-kR} below we answer Problem \ref{prob:kFU-not-kR} in the affirmative in a stronger form.

For a (zero-dimensional) Tychonoff space  $X$, we denote by $C_p(X)$ (resp., $C_p(X,\mathbf{2})$) the space $C(X)$ (resp., $C(X,\mathbf{2})$) of all continuous functions from $X$ to $\IR$ (resp., from $X$ to the doubleton $\mathbf{2}=\ZZ(2)$ considered as a discrete two element group) endowed with the pointwise topology. The space $B_1(X,\mathbf{2})$ is the space of all functions $f:X\to \mathbf{2}$ which are pointwise limits of sequences from $C_p(X,\mathbf{2})$, it is also endowed with the pointwise topology.

The characteristic function of a subset $A$ of a set $\Omega$ is denoted by $\mathbf{1}_A$.

Recall that a separable metrizable space $X$ is called a {\em $Q$-set} if every subset of $X$  is of type $F_\sigma$ in $X$. We shall use the next assertion.

\begin{proposition} \label{p:Q-set=Q2-set}
A separable zero-dimensional metrizable space $X$ is  a $Q$-set if, and only if, $B_1(X,\mathbf{2})=\mathbf{2}^X$.
%the following assertions are equivalent:
%\begin{enumerate}
%\item[{\rm(i)}] $X$ is a $Q$-set;
%\item[{\rm(ii)}] $B_1(X,\mathbf{2})=\mathbf{2}^X$.
%\end{enumerate}
\end{proposition}

\begin{proof}
Assume that $X$ is a $Q$-set. Let $f\in \mathbf{2}^X$ be arbitrary. Set  $M:=\{x\in X: f(x)=1\}$. Then $f=\mathbf{1}_M$. Since $X$ is a $Q$-set, we have $M=\bigcup_n A_n$ and $X\SM M=\bigcup_{n\in\w} B_n$ for some increasing sequences $\{A_n\}_{n\in\w}$ and  $\{B_n\}_{n\in\w}$ of closed subsets of $X$, respectively. As $X$ is zero-dimensional and Lindel\"{o}f, Theorem 6.2.7 of \cite{Eng} implies that $X$ is strongly zero-dimensional. Therefore, for every $n\in\w$, there is  a clopen subset $C_n$ of $X$ such that $A_n \subseteq C_n \subseteq X\SM B_n$. Then $\mathbf{1}_{C_n}\to \mathbf{1}_M=f$. Thus  $B_1(X,\mathbf{2})=\mathbf{2}^X$.
\smallskip

Conversely, assume that  $B_1(X,\mathbf{2})=\mathbf{2}^X$. To show that $X$ is a $Q$-set, let $M$ be an arbitrary subset of $X$. Then $\mathbf{1}_M\in B_1(X,\mathbf{2})$. Let $\{ \mathbf{1}_{A_n}\}_{n\in\w}$ be  a sequence in $C_p(X,\mathbf{2})$ which pointwise  converges to $\mathbf{1}_M$. It is clear that all $A_n$ are clopen subsets of $X$ and
\[
M=\bigcup_{k\in\w} \bigcap_{n\geq k} A_n.
\]
Since, for every $k\in\w$, the intersection $ \bigcap_{n\geq k} A_n$ is closed, we obtain that $M$ is an $F_\sigma$-set in $X$. Thus $X$ is a $Q$-set.
\end{proof}


A Tychonoff space $X$ is a {\em $k_\IR$-space} if every $k$-continuous function $f:X\to\IR$ is continuous. Recall that a function $f:X\to\IR$ is called {\em $k$-continuous} if the restriction of $f$ onto each compact subset of $X$ is continuous.
Recall also that a separable metrizable space $X$ is called a {\em $\gamma$-set} if the space $C_p(X)$ is Fr\'{e}chet--Urysohn.


\begin{example} \label{exa:kFU-count-non-kR}
Under $\mathrm{MA}$, there is a countable Boolean  $\kappa$-Fr\'{e}chet--Urysohn group $G$ which is not a $k_\IR$-space.
\end{example}

\begin{proof}
First we prove  the following claim.
\smallskip

{\em Claim 1. If there is a $\gamma$-set $X$ which is not a $Q$-set, then there is a countable Boolean  $\kappa$-Fr\'{e}chet--Urysohn group $G$ which is not a $k_\IR$-space.}
\smallskip

{\em Proof of Claim 1.} It is well-known that any $\gamma$-set is zero-dimensional, see \cite[Corollary]{GN}. It follows that the group $C_p(X,\mathbf{2})$ is separable and Hausdorff. Let $H$ be a countable and dense  subgroup of $C_p(X,\mathbf{2})$. Since $X$ is a $\gamma$-set, the space $C_p(X)$ and hence its subgroup $C_p(X,\mathbf{2})$ and the group $H$ are Fr\'{e}chet--Urysohn. As $X$  is not a $Q$-set, Proposition \ref{p:Q-set=Q2-set} implies that there is a function $g\in\mathbf{2}^X\SM B_1(X,\mathbf{2})$. Set
\[
G:=H\cup(g+H) \subseteq \mathbf{2}^X.
\]
It is clear that $G$ is a countable subgroup of $\mathbf{2}^X$. We show that $G$ is as desired.

Observe that the group $H$ is dense in $\mathbf{2}^X$ and, hence, $H$ and $g+H$ are dense in $G$. By construction, $H$ and hence also $g+H$ are Fr\'{e}chet--Urysohn spaces. Therefore, by \cite[Theorem~2.1]{Gabr-B1}, the group $G$ is a  $\kappa$-Fr\'{e}chet--Urysohn space.

It remains to show that $G$ is not a $k_\IR$-space. To this end, it suffices to show that the characteristic function $\mathbf{1}_H$ of $H$, which is discontinuous by the density of $H$ in $G$,  is $k$-continuous. To prove that $\mathbf{1}_H$ is $k$-continuous we show that for every compact subset $K$ of $G$, the sets $L:=K\cap H$ and $M:=K\cap (g+H)$ are compact. By the symmetry of $L$ and $M$, we show that $L$ is a closed subset of $G$.

Assume that $h\in \overline{L}$, and suppose for a contradiction that $h\not\in H$ and hence $h\in g+ H$.
Since $\overline{L}$ is a closed subset of the compact countable set $K$,  $\overline{L}$ is metrizable. Therefore, there is a sequence $\{g_n\}_{n\in\w}$ in $L$ converging to $h$. Take $h_0\in H$ such that $h=g+h_0$. Then $g_n-h_0\to g$, and hence $g\in B_1(X,\mathbf{2})$. But  this contradicts the choice of the function $g$ that finishes the proof of Claim 1.
\smallskip

To finish the proof of the example, we note that Theorem 1 of \cite{Gal-Miller} states that, under Martin's Axiom, there is a $\gamma$-set $X\subseteq \IR$ of cardinality the continuum. On the other hand, any $Q$-set has cardinality strictly less than $\mathfrak{c}$, see the discussion after Theorem 8.54 of \cite[p.~314]{Bukovsky2011}. Therefore the $\gamma$-set $X$ is not a $Q$-set. Now Claim 1 applies.
\end{proof}


%%%%%%%%%%%%%%%%%%
%%%%%%%%%%%%%%%%%%
%%%%%%%%%%%%%%%%%%
%%%%%%%%%%%%%%%%%%
%%%%%%%%%%%%%%%%%%





\bibliographystyle{amsplain}
\begin{thebibliography}{10}

%\bibitem{BG-Baire}
%T. Banakh, S. Gabriyelyan,  Baire category properties of some Baire type function spaces, Topology Appl. \textbf{272} (2020), 107078, 43 pp.

\bibitem{Bukovsky2011}
L. Bukovsk\'y,  {\em The structure of the real line.} Vol. 71, Springer Science \& Business Media, 2011.

\bibitem{Eng}
R.~Engelking, General Topology, Heldermann Verlag, Berlin, 1989.

\bibitem{Gabr-B1}
S. Gabriyelyan, Topological properties of spaces of Baire functions, J. Math. Anal. Appl. \textbf{478} (2019), 1085--1099.

\bibitem{Gabr-seq-Ascoli}
S. Gabriyelyan, Ascoli and sequentially Ascoli spaces, Topology Appl. \textbf{285} (2020), No 107401, 28 pp.

%\bibitem{GOR-B1}
%S. Gabriyelyan, A.V. Osipov, E. Reznichenko, Completeness and reflexivity type properties of $B_1(X)$, submitted.

\bibitem{GOR-Baire}
S. Gabriyelyan, A.V. Osipov, E. Reznichenko, Baire-type properties of topological vector spaces,  submitted.

\bibitem{GR-compact}
S. Gabriyelyan, E. Reznichenko, New classes of compact-type spaces, submitted.

\bibitem{Gal-Miller}
F. Galvin, A.W. Miller, $\gamma$-sets and other singular sets of real numbers, Topology Appl. \textbf{17} (1984), 145--155.


\bibitem{GN}
J. Gerlits, Zs. Nagy, \emph{Some properties of $C(X)$. I}, Topology Appl. \textbf{14} (1982),  151--161.

\bibitem{Ismail}
M. Ismail, Cardinal functions of homogeneous spaces and topological groups, Math. Jpn. \textbf{26} (6) (1981), 635--646.

\bibitem{kak}
J.~K\c{a}kol, W.~Kubi\'s, M.~Lopez-Pellicer, \emph{Descriptive Topology in Selected  Topics of Functional Analysis}, Developments in Mathematics, Springer, 2011.

\bibitem{LiL}
C.~Liu, L.D.~Ludwig, $\kappa$-Fr\'{e}chet--Urysohn spaces, Houston J. Math. \textbf{31} (2005), 391--401.

%\bibitem{mcoy}
%R.A. McCoy, I. Ntantu, \emph{Topological Properties of Spaces of Continuous Functions}. Lecture Notes in Math.  \textbf{1315} (1988)

\bibitem{Mrowka}
S. Mr\'{o}wka, Mazur theorem and $\mathfrak{m}$-adic spaces, Bull. Acad. Polon. Sci. S\'{e}r. Sci. Math. Astronom. Phys., 18 (1970), 299--305.

\bibitem{Osipov-23}
A.V. Osipov, The projectively Hurewicz property is $t$-invariant, Filomat \textbf{37}:28 (2023), 9613--9616.

\bibitem{Sak2}
M. Sakai, Two properties of $C_p(X)$ weaker than the Fr\'{e}echet--Urysohn property, Topol. Appl. \textbf{153}:15 (2006), 2795--2804.

\bibitem{Sakai}
M. Sakai, $\kappa$-Frechet-Urysohn property of $C_k(X)$, Topology Appl. 154 (2007) 1516--1520.

\end{thebibliography}

\end{document}
